% ----------------------
\subsection{Section 110 (3 April 1836)}
% ----------------------
	\label{DC110}
	
	% -----
	\subsubsection{Historical Background}
	% -----
		\label{DC110-HB}
		
		% --
		\paragraph{JSP}
		% --
		
			``A few days following the temple dedication in Kirtland, Ohio, and the solemn assembly that empowered church elders for the ministry, JS's journal records that JS and Oliver Cowdery had a vision of heavenly messengers in the House of the Lord. On the afternoon of Easter Sunday, 3 April 1836, JS helped other members of the church presidency distribute the sacrament of the Lord's Supper to the congregation that had assembled in the lower court of the House of the Lord. After the sacrament, the curtains were dropped, dividing the court into four quarters. According to Stephen Post, who participated in the day's meetings, the presidency then went to the pulpits for `the confirmation \& blessing of the children.' At some point during the meeting, more veils were lowered, enclosing the west pulpits and dividing them into their four levels. JS and Cowdery `retired to the pulpit'---apparently the top tier, which was reserved for the presidency---where they bowed `in solemn, but silent prayer to the Most High.'
			
			``According to the journal, after JS and Cowdery prayed, secluded in the curtains and pulpits of the temple, they had a miraculous vision of Jesus Christ, who accepted the House of the Lord as JS had prayed for at the dedication. The appearance was a fulfillment of a promise made in earlier JS revelations, that the Lord would show himself in the temple. Following the appearance of Christ, the journal records, JS and Cowdery also received visitations from the biblical prophets Moses, Elias, and Elijah, who bestowed upon the two church leaders `the Keys of this dispensation.' These keys authorized JS and Cowdery to exercise in new ways the priesthood they had received from the apostles Peter, James, and John in 1829. The bestowal of `the Keys of this dispensation,' particularly those concerning the gathering of Israel and turning `the hearts of the Fathers to the children,' marked a vital moment for Latter-day Saint missionary work and temple ordinances. Just over a year after receiving these keys, JS sent preachers to England to begin the gathering of Israel from abroad. Later, in Nauvoo, Illinois, he would teach and administer new temple ordinances that offered salvation to the deceased and bound them to the living, including baptisms for the dead, endowments, and sealings. The Latter-day Saints had shown their willingness to build the Lord a house, and these visitations on 3 April 1836 were not only a continuation of great spiritual outpouring; they were also a beginning for Latter-day Saint understanding of the purpose and power of temples.
			
			``JS and Cowdery recounted their visions to some associates shortly after they occurred. In a letter to his wife, Sally Waterman Phelps, written on the same day, William W. Phelps stated that JS and Cowdery experienced `a manifestation of the Lord' in which they learned that `the great \& terrible day of the Lord as mentioned by Malichi, was near, even at the doors.'
			
			``Sometime shortly after, Warren A. Cowdery, JS's scribe and Oliver's brother, recorded the experience in JS's journal, which is the source for the text below. Warren wrote the entry referring to JS in the third person, in contrast to the first-person language found throughout the journal. He may have relied on another original text, no longer extant, or on oral reports from either or both of the participants. If he was working from a prior text, it would directly parallel the method that produced the third-person 1834--1836 history, which he was composing in early April using JS's journal. By 7 November 1843, Willard Richards, church historian and personal secretary to JS, changed the account into first person for JS's multivolume history. JS and Oliver Cowdery's vision was added to the Doctrine and Covenants in 1876. That version, and published versions to follow, contained first-person language.
			
			``This account of visitations closes JS's 1835--1836 journal. After more than six months of almost daily recording of developments in Kirtland, entries ceased, and for nearly two years there were no entries written in this or in any other extant JS journal.''
					
	% -----
	\subsubsection{Versions}
	% -----
		\label{DC110-V}
		
		See the \href{https://drive.google.com/file/d/1goAvNz8jS4R8slYfMG_db55Ty2z_vmgK/view?usp=sharing}{sections comparison} document.
	
		\begin{compactdesc}
			\item[JSP] \hyperref[EMS-1836-JS-3-APR-1836]{3 April 1836}
			\item[BC] ---
			\item[DC 1835] ---
			\item[DC 1844] ---
			\item[TS] ---
			\item[DN] 2:26 (6 November 1852), 101
			\item[MS] 15:45--46 (5--12 November 1853), 729, 739
		\end{compactdesc}

	% -----
	\subsubsection{Section 110:1} % *DC110-1 
	% -----
		\label{DC110-1}

		THE veil was taken from our minds, and the eyes of our understanding were opened.

		\begin{framed}
		\scriptsize
			\begin{compactdesc}
				\item[JSP] The vail was taken from their minds and the eyes of their understandings were opened.
				% \item[BC] 
				% \item[]
			\end{compactdesc}
		\end{framed}

	% -----
	\subsubsection{Section 110:2} % *DC110-2 
	% -----
		\label{DC110-2}

		We saw the Lord standing upon the breastwork of the pulpit, before us; and under his feet was a paved work of pure gold, in color like amber.

		% \begin{framed}
		% \scriptsize
			% \begin{compactdesc}
				% \item[JSP] 
				% % \item[BC] 
				% % \item[]
			% \end{compactdesc}
		% \end{framed}

	% -----
	\subsubsection{Section 110:3} % *DC110-3 
	% -----
		\label{DC110-3}

		His eyes were as a flame of fire; the hair of his head was white like the pure snow; his countenance shone above the brightness of the sun; and his voice was as the sound of the rushing of great waters, even the voice of Jehovah, saying:

		% \begin{framed}
		% \scriptsize
			% \begin{compactdesc}
				% \item[JSP] 
				% % \item[BC] 
				% % \item[]
			% \end{compactdesc}
		% \end{framed}

	% -----
	\subsubsection{Section 110:4} % *DC110-4 
	% -----
		\label{DC110-4}

		I am the first and the last; I am he who liveth, I am he who was slain; I am your \hyperref[ADVOCATE]{advocate} with the Father.

		% \begin{framed}
		% \scriptsize
			% \begin{compactdesc}
				% \item[JSP] 
				% % \item[BC] 
				% % \item[]
			% \end{compactdesc}
		% \end{framed}

	% -----
	\subsubsection{Section 110:5} % *DC110-5 
	% -----
		\label{DC110-5}

		Behold, your sins are forgiven you; you are clean before me; therefore, lift up your heads and rejoice.

		% \begin{framed}
		% \scriptsize
			% \begin{compactdesc}
				% \item[JSP] 
				% % \item[BC] 
				% % \item[]
			% \end{compactdesc}
		% \end{framed}

	% -----
	\subsubsection{Section 110:6} % *DC110-6 
	% -----
		\label{DC110-6}

		Let the hearts of your brethren rejoice, and let the hearts of all my people rejoice, who have, with their might, built this house to my name.

		% \begin{framed}
		% \scriptsize
			% \begin{compactdesc}
				% \item[JSP] 
				% % \item[BC] 
				% % \item[]
			% \end{compactdesc}
		% \end{framed}

	% -----
	\subsubsection{Section 110:7} % *DC110-7 
	% -----
		\label{DC110-7}

		For behold, I have accepted this house, and my name shall be here; and I will manifest myself to my people in mercy in this house.

		% \begin{framed}
		% \scriptsize
			% \begin{compactdesc}
				% \item[JSP] 
				% % \item[BC] 
				% % \item[]
			% \end{compactdesc}
		% \end{framed}

	% -----
	\subsubsection{Section 110:8} % *DC110-8 
	% -----
		\label{DC110-8}

		Yea, I will appear unto my servants, and speak unto them with mine own voice, if my people will keep my commandments, and do not pollute this holy house.

		% \begin{framed}
		% \scriptsize
			% \begin{compactdesc}
				% \item[JSP] 
				% % \item[BC] 
				% % \item[]
			% \end{compactdesc}
		% \end{framed}

	% -----
	\subsubsection{Section 110:9} % *DC110-9 
	% -----
		\label{DC110-9}

		Yea the hearts of thousands and tens of thousands shall greatly rejoice in consequence of the blessings which shall be poured out, and the endowment with which my servants have been endowed in this house.

		% \begin{framed}
		% \scriptsize
			% \begin{compactdesc}
				% \item[JSP] 
				% % \item[BC] 
				% % \item[]
			% \end{compactdesc}
		% \end{framed}

	% -----
	\subsubsection{Section 110:10} % *DC110-10 
	% -----
		\label{DC110-10}

		And the fame of this house shall spread to foreign lands; and this is the beginning of the blessing which shall be poured out upon the heads of my people.  Even so.  Amen.

		% \begin{framed}
		% \scriptsize
			% \begin{compactdesc}
				% \item[JSP] 
				% % \item[BC] 
				% % \item[]
			% \end{compactdesc}
		% \end{framed}

	% -----
	\subsubsection{Section 110:11} % *DC110-11 
	% -----
		\label{DC110-11}

		After this vision closed, the heavens were again opened unto us; and Moses appeared before us, and committed unto us the keys of the gathering of Israel from the four parts of the earth, and the leading of the ten tribes from the land of the north.

		% \begin{framed}
		% \scriptsize
			% \begin{compactdesc}
				% \item[JSP] 
				% % \item[BC] 
				% % \item[]
			% \end{compactdesc}
		% \end{framed}

	% -----
	\subsubsection{Section 110:12} % *DC110-12 
	% -----
		\label{DC110-12}

		After this, Elias appeared, and committed the dispensation of the gospel of Abraham, saying that in us and our seed all generations after us should be blessed.

		% \begin{framed}
		% \scriptsize
			% \begin{compactdesc}
				% \item[JSP] 
				% % \item[BC] 
				% % \item[]
			% \end{compactdesc}
		% \end{framed}

	% -----
	\subsubsection{Section 110:13} % *DC110-13 
	% -----
		\label{DC110-13}

		After this vision had closed, another great and glorious vision burst upon us; for Elijah the prophet, who was taken to heaven without tasting death, stood before us, and said:

		% \begin{framed}
		% \scriptsize
			% \begin{compactdesc}
				% \item[JSP] 
				% % \item[BC] 
				% % \item[]
			% \end{compactdesc}
		% \end{framed}

	% -----
	\subsubsection{Section 110:14} % *DC110-14 
	% -----
		\label{DC110-14}

		Behold, the time has fully come, which was spoken of by the mouth of Malachi---testifying that he [Elijah] should be sent, before the great and dreadful day of the Lord come---

		% \begin{framed}
		% \scriptsize
			% \begin{compactdesc}
				% \item[JSP] 
				% % \item[BC] 
				% % \item[]
			% \end{compactdesc}
		% \end{framed}

	% -----
	\subsubsection{Section 110:15} % *DC110-15 
	% -----
		\label{DC110-15}

		To turn the hearts of the fathers to the children, and the children to the fathers, lest the whole earth be smitten with a curse---

		% \begin{framed}
		% \scriptsize
			% \begin{compactdesc}
				% \item[JSP] 
				% % \item[BC] 
				% % \item[]
			% \end{compactdesc}
		% \end{framed}

	% -----
	\subsubsection{Section 110:16} % *DC110-16 
	% -----
		\label{DC110-16}

		Therefore, the keys of this dispensation are committed into your hands; and by this ye may know that the great and dreadful day of the Lord is near, even at the doors.

		% \begin{framed}
		% \scriptsize
			% \begin{compactdesc}
				% \item[JSP] 
				% % \item[BC] 
				% % \item[]
			% \end{compactdesc}
		% \end{framed}